HiveWatch enables three classes of practitioner activity that are not well served by existing tooling.

\textbf{Portable instrumentation across FL frameworks.} Researchers who work with multiple FL frameworks — or who plan to migrate from one to another — currently must re-implement per-project monitoring code for each framework. HiveWatch provides a single instrumentation target that works across APPFL, Flower, NVIDIA FLARE, and custom training loops, reducing the instrumentation burden to five call sites regardless of which framework is used.

\textbf{Simultaneous multi-backend monitoring.} Teams frequently need to publish experiment results to more than one destination: local artifacts for immediate debugging, a hosted dashboard for cross-run comparison, and a self-hosted tracking server for compliance or reproducibility. HiveWatch's emitter model supports this with no duplication of instrumentation logic, and the SSEEmitter's replayable JSONL artifacts ensure that a run can be inspected offline after the fact without re-running training.

\textbf{FL-native diagnostics.} HiveWatch's schema captures straggler counts, gradient divergence (as a proxy for data heterogeneity), per-client staleness in asynchronous FL, communication volume breakdowns, client dropout and failure events, and geographic client metadata — none of which are available without per-project custom code when using general-purpose tracking tools such as Weights \& Biases or MLflow directly. These signals are the inputs to algorithmic heterogeneity-mitigation methods~\cite{lai2021oort,li2024fedcompass,xie2019asynchronous}, and having them logged consistently enables practitioners to characterize and diagnose heterogeneity before choosing which mitigation to apply.

The primary audience is FL practitioners and researchers who need reproducible, inspectable experiment records across diverse FL deployments — including multi-institutional biomedical FL~\cite{rieke2020future}, edge FL~\cite{lim2020federated}, and academic FL benchmarking. HiveWatch is released under the MIT license, is available on GitHub at \url{https://github.com/APPFL/hivewatch}, and is installable from source with a single \texttt{pip install} command.
